\newcommand{\etalchar}[1]{$^{#1}$}
\def\cprime{$'$} \def\Dbar{\leavevmode\lower.6ex\hbox to 0pt{\hskip-.23ex
  \accent"16\hss}D} \def\cprime{$'$}
\begin{thebibliography}{CdGH{\etalchar{+}}}

\bibitem[ACQ11]{ACQ}
Gideon Amir, Ivan Corwin, and Jeremy Quastel.
\newblock Probability distribution of the free energy of the continuum directed
  random polymer in 1 + 1 dimensions.
\newblock {\em Communications on Pure and Applied Mathematics}, 64(4):466--537,
  2011.

\bibitem[AHR98]{AHR98}
Peter~F Arndt, Thomas Heinzel, and Vladimir Rittenberg.
\newblock Spontaneous breaking of translational invariance in one-dimensional
  stationary states on a ring.
\newblock {\em Journal of Physics A: Mathematical and General}, 31(2):L45--L51,
  jan 1998.

\bibitem[Bar15]{BARRAQUAND20152674}
Guillaume Barraquand.
\newblock A phase transition for q-tasep with a few slower particles.
\newblock {\em Stochastic Processes and their Applications}, 125(7):2674--2699,
  2015.

\bibitem[BCS14]{BCSDuality}
Alexei Borodin, Ivan Corwin, and Tomohiro Sasamoto.
\newblock From duality to determinants for $q$-{TASEP} and {ASEP}.
\newblock {\em Ann. Probab.}, 42(6):2314--2382, 11 2014.

\bibitem[BDJ99]{bdj1}
Jinho Baik, Percy Deift, and Kurt Johansson.
\newblock On the distribution of the length of the longest increasing
  subsequence of random permutations.
\newblock {\em J. Amer. Math. Soc.}, 12(4):1119--1178, 1999.


\bibitem[CdGH{\etalchar{+}}]{AHR}
Zeying Chen, Jan de~Gier, Iori Hiki, Tomohiro Sasamoto, and Masato Usui.
\newblock Limiting current distribution for a two species asymmetric exclusion
  process.

\bibitem[dGMW]{GMW}
Jan de~Gier, William Mead, and Michael Wheeler.
\newblock Transition probability and total crossing events in the multi-species
  asymmetric exclusion process.

\bibitem[Dim]{Dim}
Evgeni Dimitrov.
\newblock Two-point convergence of the stochastic six-vertex model to the airy
  process.

\bibitem[DM97]{Der97}
B~Derrida and K~Mallick.
\newblock Exact diffusion constant for the one-dimensional partially asymmetric
  exclusion model.
\newblock {\em Journal of Physics A: Mathematical and General},
  30(4):1031--1046, feb 1997.

\bibitem[FV15]{FVAIHP}
Patrik~L. Ferrari and Balint Veto.
\newblock {T}racy--{W}idom asymptotics for $q$-{TASEP}.
\newblock {\em Ann. Inst. H. Poincaré Probab. Statist.}, 51(4):1465--1485, 11
  2015.



\bibitem[GO]{PatSurvey}
Patricia Goncalves and Alessandra Occelli.
\newblock On energy solutions to stochastic burgers equation.

\bibitem[Hay]{Haya}
Kohei Hayaski.
\newblock Equilibrium fluctuations for totally asymmetric interacting particle
  systems.

\bibitem[Joh00]{JohShape}
Kurt Johansson.
\newblock Shape {F}luctuations and {R}andom {M}atrices.
\newblock {\em Comm. Math. Phys}, 209:437--476, 2000.

\bibitem[KL14]{KoLe}
Marko Korhonen and Eunghyun Lee.
\newblock The transition probability and the probability for the left-most
  particle's position of the q-totally asymmetric zero range process.
\newblock {\em Journal of Mathematical Physics}, 55(1):013301, 2014.

\bibitem[KPZ86]{KPZ}
Mehran Kardar, Giorgio Parisi, and Yi-Cheng Zhang.
\newblock Dynamic scaling of growing interfaces.
\newblock {\em Physical Review Letters}, 56(9):889--892, 1986.

\bibitem[KS60]{KS60}
John~G. Kemeny and J.~Laurie Snell.
\newblock {\em Finite Markov Chains}.
\newblock Springer, 1960.

\bibitem[Kua17]{KIMRN}
Jeffrey Kuan.
\newblock A multi-species {ASEP}$(q,j)$ and $q$-{TAZRP} with stochastic
  duality.
\newblock {\em International Mathematics Research Notices},
  2018(17):5378--5416, 2017.

\bibitem[Kua18]{KuanCMP}
Jeffrey Kuan.
\newblock An algebraic construction of duality functions for the stochastic
  ${U}_q({A}_n^{(1)})$ vertex model and its degenerations.
\newblock {\em Communications in Mathematical Physics}, 359(1):121--187, Apr
  2018.

\bibitem[Kua19]{KuanAHP}
Jeffrey Kuan.
\newblock Probability distributions of multi-species q-{TAZRP} and {ASEP} as
  double cosets of parabolic subgroups.
\newblock {\em Annales Henri Poincar\'{e}}, 20(4):1149--1173, April 2019.

\bibitem[Lig76]{Ligg76}
Thomas~M. Liggett.
\newblock Coupling the simple exclusion process.
\newblock {\em Annals of Probability}, 4(3):339--356, 1976.

\bibitem[Pov13]{Pov13}
A. M. Povolotsky
\newblock On the integrability of zero-range chipping models with factorized steady states.
\newblock {\em J. Phys A}, 46(46): 465205, 2013. 

\bibitem[QS]{QS}
Jeremy Quastel and Sourav Sarkar.
\newblock Convergence of exclusion processes and KPZ equation to the KPZ fixed
  point.

\bibitem[Spi70]{Spit70}
Frank Spitzer.
\newblock Interaction of {M}arkov processes.
\newblock {\em Advances in Mathematics}, 5(2):246--290, 1970.

\bibitem[Spo14]{Spohn}
Herbert Spohn.
\newblock Nonlinear fluctuating hydrodynamics for anharmonic chains.
\newblock {\em Journal of Statistical Physics}, 154:1191–1227, 2014.

\bibitem[SW98]{SW98}
Tomohiro Sasamoto and Miki Wadati.
\newblock Exact results for one--dimensional totally asymmetric diffusion
  models.
\newblock {\em Journal of Physics A}, 31(28):6057--6071, 1998.

\bibitem[Tak]{Take15}
Yoshihiro Takeyama.
\newblock Algebraic construction of multi--species $q$--{B}oson system.

\bibitem[TW94]{TW94}
Craig~A. Tracy and Harold Widom.
\newblock Level-spacing distributions and the {A}iry kernel.
\newblock {\em Communications in Mathematical Physics}, 159(1):151--174, Jan
  1994.

\bibitem[WW09]{WW}
Jon Warren and Peter Windridge.
\newblock Some examples of dynamics for {G}elfand-{T}setlin patterns.
\newblock {\em Electron. J. Probab.}, 14:no. 59, 1745--1769, 2009.

\bibitem[WW16]{Wawa}
Dong Wang and David Waugh.
\newblock The transition probability of the q-{TAZRP} (q-bosons) with
  inhomogeneous jump rates.
\newblock {\em SIGMA}, 12(037):16, 2016.

\bibitem[Gas96]{Gasper1996ElementaryDO}
George Gasper.
\newblock Elementary derivations of summations and transformation formulas for
  q-series.
\newblock {\em arXiv: Classical Analysis and ODEs}, 1996.
\end{thebibliography}


\end{document}

@article{10.1214/aoms/1177698308,
author = {Paul R. Milch},
title = {{A Multi-dimensional Linear Growth Birth and Death Process}},
volume = {39},
journal = {The Annals of Mathematical Statistics},
number = {3},
publisher = {Institute of Mathematical Statistics},
pages = {727 -- 754},
year = {1968},
doi = {10.1214/aoms/1177698308},
URL = {https://doi.org/10.1214/aoms/1177698308}
}


To see this, 
\begin{equation}
    <\xi|\hat{S}_2S_2^*\hat{S}_1S_1^*|\eta>=\sum_{\zeta} <\xi|\hat{S}_2S_2^*|\zeta><\zeta|\hat{S}_1S_1^*|\eta>
\end{equation}
Then we have $\xi_1=\zeta_1$, $\zeta_3=\eta_3$,
and 
\begin{equation}
\begin{split}
    <\xi|\hat{S}_2S_2^*\hat{S}_1S_1^*|\eta>=D_1^{(\eta_1+\eta_2)}(\xi_1,\eta_1)q^{N(\eta_1)-N(\xi_1)}\sqrt{\mu^{(\eta_1+\eta_2)}_\alpha(\xi_1)\mu^{(\eta_1+\eta_2)}_\alpha(\eta_1)}(-1)^{N(\eta_1)}\\
  \times D_2^{(\theta-\xi_1)}(\xi_2,\eta_1+\eta_2-\xi_1)q^{N(\eta_1+\eta_2-\xi_1)-N(\xi_2)}\\
 \times  \sqrt{\mu^{(2j-\xi_1)}_\alpha(\xi_2)\mu^{(2j-\xi_1)}_\alpha(\eta_1+\eta_2-\xi_1)}(-1)^{N(\eta_1+\eta_2-\xi_1)}\\
\end{split}
\end{equation}

%\begin{equation}
 %   \langle\xi|F_2F_1|\eta\rangle=\sum_{\zeta} \langle\xi|F_2|\zeta\rangle\langle\zeta|F_1|\eta\rangle
%\end{equation}


\begin{equation}
\begin{split}
  \,_{N}  \langle\xi|F_2F_1|\eta\rangle_{N}=D_1^{(\eta_1+\eta_2)}(\xi_1,\eta_1)\sqrt{\mu^{(\eta_1+\eta_2)}_\alpha(\xi_1)\mu^{(\eta_1+\eta_2)}_\alpha(\eta_1)}\\
  \times q^{\binom{N(\xi_1)}{2}-\binom{N(\eta_1)}{2}}\frac{\left(\alpha q^{1+2N(\xi_1)-2N(\eta_1+\eta_2)} \right)_\infty}{\left(\alpha q^{1-2N(\eta_1)} \right)_\infty} \\
  \times D_2^{(N-\xi_1)}(\xi_2,\eta_1+\eta_2-\xi_1)
 \sqrt{\mu^{(N-\xi_1)}_\alpha(\xi_2)\mu^{(N-\xi_1)}_\alpha(\eta_1+\eta_2-\xi_1)}\\ \times q^{\binom{N(\xi_2)}{2}-\binom{N(\eta_1+\eta_2-\xi_1)}{2}}\frac{\left(\alpha q^{1+2N(\xi_2)-2N(N-\xi_1)} \right)_\infty}{\left(\alpha q^{1-2N(\eta_1+\eta_2-\xi_1)} \right)_\infty} \\
\end{split}
\end{equation}




Then $\frac{\,_{N}  \langle\xi|F_2F_1|\eta\rangle_{N}}{\sqrt{ \mu_\alpha^3(\xi) \mu_\alpha^3(\eta)}}$ is orthogonal with respect to measure $ \mu_\alpha^3$, and the duality function is 

\begin{equation}
  \mathcal{D}=  \frac{\,_{N}  \langle\xi|F_2F_1|\eta\rangle_{N}}{\sqrt{ \mu_\alpha^3(\xi) \mu_\alpha^3(\eta)}}
\end{equation}


These values are expressed as double integrals of Laplace type, for which a general formula for the asymptotic \textbf{expansion} from Chapter 8 of \cite{BLBook} will be cited here.

\begin{theorem}
Define 
$$
I(\lambda) = \int_D \exp(\lambda \phi(\mathbf{x})) g_0(\mathbf{x})d \mathbf{x}
$$
where $D$ is a domain in $\mathbb{R}^n$ and that $\phi,g_0$ are smooth functions in $D$. Further suppose that $\phi$ obtains its absolute maximum at some point $\mathbf{x}_0 \in D$. Then for any positive integer $M$,
$$
I(\lambda) \sim \exp(\lambda\phi(\mathbf{x}_0)) \left( \frac{2\pi}{\lambda}\right)^{n/2} \sum_{j=0}^{M-1} \frac{\Delta_{\xi}^j G_0 \vert_{\xi=0}}{(2\lambda)^jj!},
$$
where $G_0(\mathbf{\xi})$ is the function defined as follows. Let $A$ be the $n\times n$ matrix with entries
$$
A_{ij} = \phi_{x_ix_j}(\mathbf{x}_0),
$$
and let $Q$ be an orthogonal matrix with diagonalizes $A$ in such a way that
$$
Q^TAQ =: \Lambda = \mathrm{diag}(\lambda_1,\ldots,\lambda_n)
$$
where $\lambda_1,\ldots,\lambda_n$ are the eigenvalues of $A$. Now set
$$
(\mathbf{x} - \mathbf{x}_0)^T = QR \mathbf{z}^T, \quad R = \mathrm{diag}(| \lambda_1|^{-1/2}, \ldots, | \lambda_n|^{-1/2})
$$
and 
$$
f(\mathbf{z}) := \phi(\mathbf{x}_0) - \phi(\mathbf{x}(\mathbf{z})).
$$
Introduce the (non--unique) quantities $\xi_i = h_i(\mathbf{z}), i=1,2,\ldots,n$ such that $$h_i = z_i + o(\vert \mathbf{z} \vert) \text{ as }\vert \mathbf{z} \vert \rightarrow 0$$ and
$$
\sum_{i=1}^n h_i^2 = 2f.
$$
Then
$$
G_0(\xi) = g_0(\mathbf{x}(\mathbf{\xi}))J(\xi)
$$
where $J(\xi) = \partial(x_1,\ldots,x_n)/\partial(\xi_1,\ldots,\xi_n)$ is the Jacobian. The asymptotic expansion does not depend on the choice of $h_i$. 
\end{theorem}
\begin{proof}
This is equation (8.3.50).
\end{proof}


Furthermore, replacing $w_j = \omega_j/L$ we obtain
\begin{multline*}
\mathcal{E}+\frac{q^{N(N-1)/2}L^{N/2}}{(2\pi i )^N} \int_{-\infty}^{\sigma_1 } d\nu_1 \cdots \int_{-\infty}^{\sigma_N} d\nu_N\\
\times \int \frac{d\omega_1}{\omega_1} \cdots \int \frac{d\omega_N}{\omega_N} B(\omega_1,\ldots,\omega_N) \prod_{j=1}^N \log(1-\frac{\omega_j}{L}) (1-\frac{\omega_j}{L})^{-L-L^{1/2}\nu_j} e^{-\omega_j  }   ,
\end{multline*}
where $\mathcal{E}$ is an error term obtained from replacing $(M_j-L)L^{-1/2}$ with $\sigma_j$. Now recall that
$$
e^{-\omega}\left( 1 - \frac{\omega}{L}\right)^{-L} = 1 +  \frac{\omega^2}{2L} + O \left(\frac{\omega^3}{2L^2}\right)
$$
Let $\mathcal{E}'$ be the second error term that occurs from the $O(\cdot)$ approximation. Then switching back to $w_j$, 
\begin{multline*}
\mathcal{E}+ \mathcal{E}'+\frac{q^{N(N-1)/2}L^{N/2}}{(2\pi i )^N} \int_{-\infty}^{\sigma_1 } d\nu_1 \cdots \int_{-\infty}^{\sigma_N} d\nu_N\\
\times \int dw_1 \cdots \int dw_N B(w_1,\ldots,w_N) \prod_{j=1}^N \frac{\log(1-w_j)}{w_j} (1-{w_j})^{-L^{1/2}\nu_j} \left( 1 + \frac{w_jL}{2} \right).
\end{multline*}
Now recall that
$$
\frac{\log(1-w)}{w} = -1 - \frac{w}{2} - \frac{w^2}{3} - \frac{w^3}{4} - \cdots 
$$
so that
\begin{multline*}
\mathcal{E}+ \mathcal{E}'+\frac{q^{N(N-1)/2}L^{N/2}}{(2\pi i )^N} \int_{-\infty}^{\sigma_1 } d\nu_1 \cdots \int_{-\infty}^{\sigma_N} d\nu_N\\
\times \int dw_1 \cdots \int dw_N B(w_1,\ldots,w_N) \prod_{j=1}^N (-1-w_j+O(w_j^2))  e^{(w_jL^{1/2}+w_j^2L^{1/2}/2+O(w_j^3))\nu_j} .
\end{multline*} 
Now substitute $w_j = iu_jL^{-1/2}$, so that 




\textcolor{orange}{maybe we should delete the 2-species case? YES. Will DO.}

Recall the reversible measures for ASEPs  are \textcolor{orange}{It should also be deleted? appeared before YES TOO.}
\begin{equation}
    \mu_{\mathbf{k}}^n(\bs\xi)=1_{\{N(\bs\xi_i)=k_i,\ 0\le i\le n\}} \prod_x[\theta^x]_q!\prod_{i=0}^{n}\frac{1}{[\xi_i^x]^!_q}q^{(\xi_i^x)^2/2}\prod_{y<x}\prod_{i=0}^{n-1}q^{-2\xi^x_{[0,i]}\xi^y_{i+1}},
\end{equation}
\begin{equation}
    \mu_\alpha^{\bs\theta}(\bs\xi)=\prod_{x=1}^L \alpha^{\xi^x}\binom{\theta^x}{\xi^x}_qq^{-\left(2N^-_{x-1}(\bs\theta)+\theta^x\right)\xi^x}\\
    =q^{-2N(\bs\theta)N(\bs\xi)}\prod_{x=1}^L \alpha^{\xi^x}\binom{\theta^x}{\xi^x}_qq^{\left(2N^+_{x+1}(\bs\theta)+\theta^x\right)\xi^x}.
\end{equation}

\begin{theorem}
\begin{equation}
   q^{h(\bs\xi,\bs\eta)} \prod_{x=1}^L{}_1\phi_0\left(q^{-2\xi_0^x};q^2,q^{-2\left(N_{x-1}^-(\bs\xi_0)+N_{x+1}^+(\bs\eta_1)\right)+1}\right)
    \prod_{x=1}^L{}_1\phi_0\left(q^{-2\xi_1^x};q^2,q^{-2\left(N_{x-1}^-(\bs\xi_1)+N_{x+1}^+(\bs\eta_0)\right)+1}\right).
\end{equation}
is a space reversed duality for $2$-species $q$-TAZRP, where $\bs\eta$ evolves with total asymmetry to the left and $\bs\xi$ evolves with total asymmetry to the right. 
\begin{equation}
  h(\bs\xi,\bs\eta)%= -\sum_x\xi_1^xN_{x+1}^+(\eta_1)+\sum_x\eta_1^xN_{x+1}^+(\xi_1)+\sum_x\xi_1^x\eta_1^x\\
  = -\sum_x\xi_0^xN_{x+1}^+(\bs\eta_0)+\sum_x\eta_0^xN_{x}^+(\bs\xi_0).
\end{equation}
\end{theorem}

\begin{proof}



For $2$-species ASEP$(q,\bs\theta)$, we investigate the duality function
\begin{equation}
  \mathcal{D}(\bs\xi,\bs\eta)= D_{\alpha_0}^{(\bs\eta_0+\bs\eta_1)}(\bs\xi_0,\bs\eta_0) D_{\alpha_1}^{(\bs\theta-\bs\xi_0)}(\bs\xi_1,\bs\eta_0+\bs\eta_1-\bs\xi_0)
  \sqrt{\frac{\mu^{(\bs\eta_0+\bs\eta_1)}_{\alpha_0}(\xi_0)\mu^{(\bs\eta_0+\bs\eta_1)}_{\alpha_0}(\bs\eta_0)\mu^{(\bs\theta-\bs\xi_0)}_{\alpha_1}(\bs\xi_1)\mu^{(\bs\theta-\bs\xi_0)}_{\alpha_1}(\bs\eta_0+\bs\eta_1-\bs\xi_0)}{\mu^3(\bs\xi) \mu^3(\bs\eta)}},
\end{equation}
\begin{equation}
  \mathcal{D}(\bs\xi,T(\bs\eta))= D_{\alpha_0}^{(\bs\theta-\bs\eta_0)}(\bs\xi_0,\bs\theta-\bs\eta_0-\bs\eta_1) D_{\alpha_1}^{(\bs\theta-\bs\xi_0)}(\bs\xi_1,\bs\theta-\eta_0-\xi_0)\sqrt{\frac{\mu^{(\bs\theta-\bs\eta_0)}_{\alpha_0}(\bs\xi_0)\mu^{(\bs\theta-\eta_0)}_{\alpha_0}(\bs\theta-\bs\eta_0-\bs\eta_1)\mu^{(\bs\theta-\bs\xi_0)}_{\alpha_1}(\bs\xi_1)\mu^{(\bs\theta-\bs\xi_0)}_{\alpha_1}(\bs\theta-\bs\eta_0-\bs\xi_0)}{\mu^3(\bs\xi) \mu^3(T(\bs\eta))}}
\end{equation}
Take $\alpha_0=q^{2N(\bs\theta-\bs\eta_0)}$, $\alpha_1=q^{2N(\bs\theta-\bs\xi_0)}$,
\begin{equation}
  \lim_{\bs\theta\xrightarrow{}\infty}  D_{\alpha_0}^{(\bs\theta-\bs\eta_0)}(\bs\xi_0,\bs\theta-\bs\eta_0-\bs\eta_1) =\prod_{x=1}^L{}_1\phi_0\left(q^{-2\xi_0^x};q^2,q^{-2\left(N_{x-1}^-(\bs\xi_0)+N_{x+1}^+(\bs\eta_1)\right)+1}\right)
\end{equation}
\begin{equation}
  \lim_{\bs\theta\xrightarrow{}\infty}  D_{\alpha_1}^{(\bs\theta-\xi_0)}(\bs\xi_1,\bs\theta-\bs\eta_0-\bs\xi_0) =\prod_{x=1}^L{}_1\phi_0\left(q^{-2\xi_1^x};q^2,q^{-2\left(N_{x-1}^-(\bs\xi_1)+N_{x+1}^+(\bs\eta_0)\right)+1}\right)
\end{equation}
\begin{equation}
\lim_{\theta^x\xrightarrow{}\infty} \sqrt{ \frac{\binom{\theta^x-\eta_0^x}{\xi_0^x}_q     \binom{\theta^x-\eta_0^x}{\theta^x-\eta_0^x-\eta_1^x}_q    \binom{\theta^x-\xi_0^x}{\xi_0^x}_q       \binom{\theta^x-\xi_0^x}{\theta^x-\eta_0^x-\xi_0^x}_q}{ \binom{\theta^x}{\xi^x}_q  \binom{\theta^x}{T(\bs\eta)^x}_q }} =\lim _{\theta^x \rightarrow \infty} \frac{\left(\begin{array}{c}\theta^x-\eta_{0}^{x} \\ \xi_0^{x}\end{array}\right)_{q}}{\left(\begin{array}{l}\theta^x \\ \xi_{0}^{x}\end{array}\right)_{q}}=q^{\xi_{0}^{x} \eta_{0}^{x}}
\end{equation}
Last, investigate the power of $q$  in $\mathcal{D}(\bs\xi,T(\bs\eta))$ that comes from the reversible measures.
Compute it explicitly, up to a constant, we have %the measure is the second one of (1.8)

\begin{multline}
  r(\bs\theta,\bs\xi,\bs\eta) =    \sum_x (2N_{x+1}^+(\bs\theta-\bs\eta_0)+\theta^x-\eta_0^x)(\theta^x+\xi_0^x-\eta_0^x-\eta_1^x)
        +\sum_x (2N_{x+1}^+(\bs\theta-\bs\xi_0)+\theta^x-\xi_0^x)(\theta^x+\xi_1^x-\eta_0^x-\xi_0^x)\\
    -\sum_x\sum_{i=0}^2\frac{(\xi_i^x)^2+(\eta_i^x)^2}{2}
        +\sum_x\sum_{y<x}\sum_{i=0}^12\xi^x_{[0,i]}\xi_{i+1}^y +\sum_x\sum_{y<x}\sum_{i=0}^12T(\bs\eta)^x_{[0,i]}T(\bs\eta)_{i+1}^y\\ -2N(\bs\theta-\bs\eta_0)N(\bs\xi)-2N(\bs\theta-\bs\eta_0)N(\bs\theta-\bs\eta_0-\bs\eta_1)
     -2N(\bs\theta-\bs\xi_0)N(\bs\xi_1)-2N(\bs\theta-\bs\xi_0)N(\bs\theta-\bs\eta_0-\bs\xi_0)   .   \end{multline}


If this term can be written as  sum of a function $f(\bs\theta,N(\bs\xi),N(\bs\eta))$ and a function $g(\bs\xi,\bs\eta)$ that does not depend on $\bs\theta$, we can divide the duality function by $q^{f/2}$, thus get a nontrivial limit 
\begin{multline}
  \lim_{\bs\theta\xrightarrow{}\infty} \frac{\mathcal{D}(\bs\xi,T(\bs\eta)) }{q^{\frac{f(\bs\theta,N(\bs\xi),N(\bs\eta))}{2}}} = q^{\frac{g(\bs\xi,\bs\eta)}{2}+\sum_x\xi_0^x\eta_0^x} \prod_{x=1}^L{}_1\phi_0\left(q^{-2\xi_0^x};q^2,q^{-2\left(N_{x-1}^-(\bs\xi_0)+N_{x+1}^+(\bs\eta_1)\right)+1}\right)\\
   \times \prod_{x=1}^L{}_1\phi_0\left(q^{-2\xi_1^x};q^2,q^{-2\left(N_{x-1}^-(\bs\xi_1)+N_{x+1}^+(\bs\eta_0)\right)+1}\right).  
\end{multline}


Computing $f,g$, we have 
\begin{equation}\label{eq:2}
    \begin{split}
        g(\bs\xi,\bs\eta)%=-\sum_x(2N_{x+1}^+(\eta_1)+\eta_1^x)(\xi_1^x-\eta_1^x-\eta_2^x)-\sum_x(2N_{x+1}^+(\xi_1)+\xi_1^x)(\xi_2^x-\xi_1^x-\eta_1^x)\\
        %-\sum_x\frac{(\xi_1^x)^2+(\xi_2^x)^2+(\xi_1^x+\xi_2^x)^2+(\eta_1^x)^2+(\eta_2^x)^2+(\eta_1^x+\eta_2^x)^2}{2}\\
       % +\sum_x\sum_{y<x}(2\xi_1^x\xi_2^y-2(\xi_1^x+\xi_2^x)(\xi_1^y+\xi_2^y))-\sum_x\sum_{y<x}(2\eta_1^x\eta_1^y+2(\eta_1^x+\eta_2^x)\eta_2^y)\\
       % +2N(\xi_1)N(\xi_2)+N(\xi_1)^2+N(\eta_2)^2\\
        =-2\sum_x\xi_0^xN_{x+1}^+(\bs\eta_0)+2\sum_x\eta_0^xN_{x+1}^+(\bs\xi_0).
    \end{split}
\end{equation}
\end{proof}

